\documentclass{article}
\usepackage[final]{neurips_2022}
\usepackage[utf8]{inputenc}
\usepackage[T1]{fontenc}
\usepackage[colorlinks=true,bookmarksnumbered,
            citecolor=darkblue,urlcolor=darkblue,linkcolor=darkblue]{hyperref}
\usepackage{url,booktabs,amsfonts,nicefrac,microtype,xcolor}
\usepackage{natbib}
\usepackage{amsmath,amssymb,amsthm}
\usepackage{graphicx,xspace}
\definecolor{darkblue}{rgb}{0,0,0.75}

\newcommand{\GNNEx}{\textsc{GNNExplainer}\xspace}
\newcommand{\MI}{\mathrm{MI}}
\newcommand{\reals}{\mathbb{R}}
\newcommand{\calG}{\mathcal{G}}
\newcommand{\calV}{\mathcal{V}}
\newcommand{\calE}{\mathcal{E}}
\newcommand{\calN}{\mathcal{N}}
\newcommand{\calY}{\mathcal{Y}}
\newcommand{\bfA}{\mathbf{A}}
\newcommand{\bfX}{\mathbf{X}}
\newcommand{\bfH}{\mathbf{H}}
\newcommand{\bfM}{\mathbf{M}}
\newcommand{\bfm}{\mathbf{m}}
\newcommand{\bfx}{\mathbf{x}}
\newcommand{\bfz}{\mathbf{z}}
\newcommand{\bfw}{\mathbf{w}}
\newcommand{\sigmoid}{\sigma}
\DeclareMathOperator*{\argmax}{arg\,max}

\theoremstyle{definition}
\newtheorem{problem}{Problem}

\title{GNNExplainer: Generating Explanations for Graph Neural Networks}

\author{%
  Arion Scheid\\
  Department of Computer Science\\
  TU Wien\\
  Vienna, Austria\\
} 

\begin{document}
\maketitle

\begin{abstract}
Graph Neural Networks (GNNs) achieve strong performance on graph-structured data
but are largely opaque: it is unclear which nodes, edges, or features drive a
given prediction. This paper examines \GNNEx \citep{ying2019gnnexplainer}, a
post-hoc method that finds the compact subgraph and feature subset that maximize
mutual information with a trained model's prediction. We present a
self-contained derivation of this objective and its continuous relaxation into
learned edge and feature masks, including a step we show is mis-stated in the
original derivation, connecting it to the practical optimization implemented
in PyTorch Geometric \citep{fey2019fast}. We study two failure modes. First,
gradient-based mask optimization can become trapped in local optima; on a
novel multi-size clique dataset, explanation F1 drops from 0.12 for triangles
to 0.00 for 5-cliques. Second, motivated by 2026 work showing that
jointly-trained \emph{self-explainable} GNNs can produce explanations
provably unrelated to a model's actual reasoning
\citep{azzolin2026gnnexplanations}, we add two constant, label-independent
color features to the clique dataset and test whether \GNNEx's feature
mask mistakes either one for a real signal. We find only a weak and
inconsistent version of this failure: the mask does not fully suppress
either feature, but the class-dependent gap is small and does not track
which class the classifier finds harder to predict.
\end{abstract}


\section{Introduction}

Graph Neural Networks have become the standard approach for learning on
graph-structured data, with applications in molecular property prediction,
fraud detection, and scientific knowledge graphs. Despite their success, GNNs
are black boxes: a prediction cannot easily be traced to specific input
elements. This is not merely an aesthetic concern, since a predicted
toxicity must be linked to a functional group to guide synthesis, and a
flagged transaction must be auditable for regulators.

The general problem of explaining black-box predictions has been studied
extensively \citep{lipton2018mythos}, but standard tools such as LIME
\citep{ribeiro2016should} and SHAP \citep{lundberg2017unified} treat inputs as
flat feature vectors and cannot account for graphs' relational structure, so a
dedicated method is needed.

\GNNEx \citep{ying2019gnnexplainer} fills this gap. Given any trained GNN and a
node to explain, it identifies the minimal subgraph and node feature subset
that best account for the prediction, by maximizing \emph{mutual information}
with the model's output: a good explanation, fed to the frozen GNN alone,
reproduces the original prediction as confidently as possible. \GNNEx
operationalizes this with a continuous, differentiable \emph{mask} over the
edges and features of the local computation graph, optimized by gradient
descent. The method is post-hoc and model-agnostic: it works on any GNN
without retraining, and the GNN's parameters stay \emph{frozen} throughout.

This paper makes four contributions. First, we derive the mutual information
objective from scratch, identifying a step in the original derivation where
Jensen's inequality is applied in the wrong direction
(Section~\ref{subsec:jensen}). Second, we concretely investigate the
local-optima behavior that \citet{ying2019gnnexplainer} already anticipate
but do not characterize: when it settles on sub-cliques or sub-motifs that
are locally but not globally optimal. Third, we connect \GNNEx to 2026 work
on \emph{self-explainable} GNNs showing that some explanation mechanisms
can be provably unrelated to a model's actual reasoning, adapting this
construction into a testable hypothesis for \GNNEx. Fourth, we implement
\GNNEx on a multi-size clique benchmark with known ground-truth
explanations, demonstrating the local-optima behavior and empirically
testing whether the unfaithfulness failure above also affects \GNNEx's
feature mask.


\section{Related Work}
\label{sec:related}

\paragraph{Graph Neural Networks.}
The message-passing framework of \citet{gilmer2017neural} unifies most GNN
architectures. Examples of such are GCN (Graph Convolutional Networks) of \citet{kipf2017semi} which use symmetric normalized
aggregation or GAT (Graph Attention Networks) of \citet{velickovic2018graph} which use learned
attention weights, sometimes interpreted as explanations (as in \cite{ying2019gnnexplainer}), though
\citet{jain2019attention} show high attention on an edge does not imply that
edge is necessary for the prediction.

\paragraph{GNN Explanation Methods.}
A comprehensive taxonomy of GNN explanation methods is provided by
\citet{yuan2022explainability}. Among instance-level methods, \GNNEx
\citep{ying2019gnnexplainer} is the seminal, most widely cited method and
the primary reference for this paper. \citet{li2020explain}
extend the idea to weighted graphs. PGExplainer \citep{luo2020parameterized}
addresses \GNNEx's consistency limitation with a shared generative model
instead of a separate optimization per instance. SubgraphX
\citep{yuan2021explainability} avoids gradient-based optimization entirely,
searching combinatorially instead. XGNN
\citep{yuan2020xgnn} trains a graph generator at the model level.

\paragraph{Evaluation of GNN Explanations.}
A substantial literature asks not how to produce GNN explanations but
whether to trust them. \citet{faber2021when}, on evaluation methodology,
identify five pitfalls in comparing explanations to planted ground truth
(Section~\ref{sec:limitations}), and \citet{agarwal2022probing} give the
first theoretical reliability analysis, bounding faithfulness, stability,
and fairness. This report sits in the same field: Section~\ref{sec:results}
applies both papers' cautions to our results, and Section~\ref{subsec:anchor}
extends the same question to a construction from
\citet{azzolin2026gnnexplanations}, discussed next.

\paragraph{Self-Explainable GNNs and Unfaithful Explanations.}
\citet{azzolin2026gnnexplanations} study this question for
\emph{self-explainable} GNNs (SE-GNNs). An SE-GNN couples an
\emph{explanation extractor} $e$, mapping an input graph to a subgraph
$e(G) = R \subseteq G$, with a \emph{classifier} $g$ that predicts from $R$
alone, so the full model is $f(G) = g(e(G))$; GSAT \citep{miao2022gsat} is a
representative instance. Since $e$ and $g$ train jointly, $e$ can discard the
information $g$ needs and instead report a recurrent, non-discriminative
pattern: an \emph{anchor set} $Z = \{z_y\}_{y \in \calY}$ of single-node
subgraphs, one per class, with $z_y \subseteq G$ in \emph{every} instance $G$
(Appendix~\ref{app:rbgv} gives a worked example). High accuracy therefore
cannot certify that an SE-GNN's explanation reflects its real computation.
Section~\ref{subsec:anchor} asks whether a weaker analogue of this result
survives for post-hoc explainers such as \GNNEx.


\section{Background: Graph Neural Networks}
\label{sec:background}

A graph is a tuple $\calG = (\calV, \calE, \bfX)$ with node set $\calV$ of
size $n = |\calV|$, edge set $\calE$, and node feature matrix
$\bfX \in \reals^{n \times d}$, where $d$ is the number of features per node.
The adjacency matrix $\bfA \in \{0,1\}^{n \times n}$ encodes the edge set. We
assume undirected graphs throughout and write
$\calN(v) = \{u \in \calV : (v,u) \in \calE\}$ for the neighbors of node $v$.

GNNs are trained for a classification objective: given $\calG$, predict a
label from a fixed set of $C$ classes $\{1, \dots, C\}$. This prediction can
be made at several levels. Node classification predicts a true label
$y_v \in \{1, \dots, C\}$ for every node $v \in \calV$. Link prediction
predicts a label for every pair of nodes. Graph classification predicts a
single label for the whole graph. We focus on node classification
throughout, following \citet{kipf2017semi}.

A $K$-layer GNN, which we denote $\Phi$, computes node embeddings by
iterating the message-passing scheme of \citet{gilmer2017neural}. At layer
$k$, each node $v$ first aggregates information from its neighbors and then
updates its own representation:
\begin{align}
  \bfm_v^{(k)} &= \mathrm{AGG}^{(k)}\!\bigl(
    \{\phi^{(k)}(\bfH_v^{(k-1)}, \bfH_u^{(k-1)}, \bfA_{vu})
      : u \in \calN(v)\}\bigr), \label{eq:agg} \\
  \bfH_v^{(k)} &= \mathrm{UPDATE}^{(k)}\!\bigl(\bfH_v^{(k-1)},
                                                 \bfm_v^{(k)}\bigr),
  \label{eq:update}
\end{align}
with initialization $\bfH_v^{(0)} = \bfX_v$, the raw node features. Here
$\bfH_v^{(k)} \in \reals^{d_k}$ is the \emph{hidden representation}
(embedding) of node $v$ after $k$ layers. $\bfm_v^{(k)}$ is the
\emph{aggregated message} received from $v$'s neighbors at layer $k$.
$\phi^{(k)}$ is a learnable function, for example a linear transformation,
that combines a pair of node embeddings with the edge weight $\bfA_{vu}$.
$\mathrm{AGG}^{(k)}$ is a permutation-invariant aggregation function over the
neighbor set, for example a sum or mean. $\mathrm{UPDATE}^{(k)}$ combines the
current node embedding $\bfH_v^{(k-1)}$ with the incoming message
$\bfm_v^{(k)}$, typically with a linear layer followed by a non-linearity.

The final embedding $\bfH_v^{(K)}$ is passed through a softmax layer that
outputs a probability $P_\Phi(Y\!=\!c \mid \calG, \bfX)$ for each class $c$,
where $Y$ is the random variable representing the label $\Phi$ predicts for
$v$. The predicted label is the most probable class,
\begin{equation}
  \hat{y}_v = \argmax_{c \in \{1,\dots,C\}} P_\Phi(Y\!=\!c \mid \calG, \bfX).
  \label{eq:pred}
\end{equation}

A critical structural property is that $\hat{y}_v$ depends only on the
$K$-hop neighborhood of $v$. We call this neighborhood the
\emph{computation graph} $\calG_c^v$, with associated features $\bfX_c^v$.
All nodes and edges outside $\calG_c^v$ are irrelevant to $\hat{y}_v$, so
$P_\Phi(Y\!=\!c \mid \calG, \bfX) = P_\Phi(Y\!=\!c \mid \calG_c^v, \bfX_c^v)$:
only $\calG_c^v$ matters for explaining $\hat{y}_v$.


\section{GNNExplainer}
\label{sec:method}

\subsection{Problem Formulation}

The GNN $\Phi$ of Section~\ref{sec:background} is trained once and then kept
\emph{frozen} during explanation. The goal is to identify the part of the
computation graph $\calG_c^v$ that best explains the prediction $\hat{y}_v$.

\begin{problem}[Instance-Level GNN Explanation \citep{ying2019gnnexplainer}]
Given $\Phi$, a computation graph $\calG_c^v$ with features $\bfX_c^v$, and
prediction $\hat{y}_v$, find a subgraph $\calG_S \subseteq \calG_c^v$ and a
masked feature matrix $\bfX_S^F$ such that $(\calG_S, \bfX_S^F)$ is most
\emph{informative} about $\hat{y}_v$, i.e., knowing $(\calG_S, \bfX_S^F)$
reduces uncertainty about $\hat{y}_v$ as much as possible.
\end{problem}

\subsection{Information-Theoretic Preliminaries}

\GNNEx formalizes ``most informative'' using \emph{mutual information}. We
first define the necessary information-theoretic quantities for the random
variable $Y$ of Eq.~\eqref{eq:pred}, which ranges over the $C$ classes
$\{1,\dots,C\}$ of Section~\ref{sec:background}.

The \emph{Shannon entropy} measures the uncertainty of~$Y$:
\begin{equation}
  H(Y) = -\sum_{c=1}^{C} P(Y\!=\!c)\,\log P(Y\!=\!c).
  \label{eq:entropy}
\end{equation}
$H(Y) \geq 0$, with equality only when one class has probability~1 (no
uncertainty), and is maximized when all classes are equally likely.
The \emph{conditional entropy} measures the residual uncertainty in~$Y$
after observing evidence~$Z$:
\begin{equation}
  H(Y \mid Z) = -\sum_{c=1}^{C} P(Y\!=\!c \mid Z)\,\log P(Y\!=\!c \mid Z),
  \label{eq:cond_entropy_def}
\end{equation}
with $H(Y \mid Z) \leq H(Y)$: conditioning on informative evidence can
only reduce uncertainty.  \emph{Mutual information} quantifies this
reduction:
\begin{equation}
  \MI(Y;\, Z) = H(Y) - H(Y \mid Z) \;\geq\; 0.
  \label{eq:mi_def}
\end{equation}
A high value of $\MI(Y; Z)$ means observing $Z$ makes $Y$ substantially
more predictable.  For example, a confident 3-class prediction $P =
(0.9, 0.05, 0.05)$ has $H \approx 0.47$\,bits, while a near-uniform
prediction $P = (0.33, 0.33, 0.34)$ has $H \approx 1.10$\,bits.

\subsection{Mutual Information Objective}

\GNNEx seeks the explanation $(\calG_S, \bfX_S^F)$ that maximizes the mutual
information between $Y$ and this evidence:
\begin{equation}
  \max_{\calG_S \subseteq \calG_c^v}\ \MI\bigl(Y;\, (\calG_S, \bfX_S^F)\bigr)
  = \underbrace{H(Y)}_{\text{constant}} - H\bigl(Y \mid \calG\!=\!\calG_S,\,
                                       \bfX\!=\!\bfX_S^F\bigr).
  \label{eq:mi}
\end{equation}
Since $\Phi$ is fixed, $H(Y)$ is a constant.  Maximizing MI is therefore
equivalent to \emph{minimizing the conditional entropy}, that is, finding the
subgraph that keeps the GNN most confident:
\begin{equation}
  \min_{\calG_S,\, \bfX_S^F}\ H\bigl(Y \mid \calG\!=\!\calG_S,\,
                                       \bfX\!=\!\bfX_S^F\bigr)
  = -\sum_{c=1}^{C} P_\Phi(Y\!=\!c \mid \calG_S, \bfX_S^F)\,
    \log P_\Phi(Y\!=\!c \mid \calG_S, \bfX_S^F).
  \label{eq:cond_ent}
\end{equation}
This uses the \textbf{entire softmax distribution} over all $C$ classes and
penalizes any residual uncertainty.

\paragraph{Class-specific formulation.}
In practice, the GNN's predicted label $\hat{y}_v = c^*$ is known.
\citet{ying2019gnnexplainer} replace the full conditional entropy with a
\emph{cross-entropy targeted at $c^*$}:
\begin{equation}
  \min_{\calG_S,\, \bfX_S^F}\;
  -\log P_\Phi\bigl(Y\!=\!c^* \mid \calG\!=\!\calG_S,\,
  \bfX\!=\!\bfX_S^F\bigr).
  \label{eq:class_specific}
\end{equation}
While Eq.~\eqref{eq:cond_ent} asks \emph{``Is the GNN confident?''} across all
classes, Eq.~\eqref{eq:class_specific} asks \emph{``Does $\calG_S$ support
predicting~$c^*$?''} by only looking at $P(Y\!=\!c^*)$. The two agree when
the GNN is highly confident. Eq.~\eqref{eq:class_specific} is simpler,
directly interpretable as the standard cross-entropy loss used in
classification, and is the form implemented in PyTorch Geometric.

\subsection{Mask Parameterization}
\label{subsec:maskparam}

Searching over all $2^{|\calE_c^v|}$ subgraphs by brute force is intractable.
\GNNEx sidesteps this with continuous \emph{masks}: real-valued scalars (one
per edge or feature dimension) passed through a sigmoid to produce soft
weights in $(0,1)$, optimized jointly by gradient descent, with the final
discrete subgraph recovered by thresholding.

\paragraph{Edge mask.}
An \emph{edge mask} is a real-valued matrix $\bfM_A \in \reals^{n \times n}$
learned with one entry per edge in $\calG_c^v$. Applying the sigmoid $\sigmoid(\cdot)$
gives a soft weight in $(0,1)$, and the masked adjacency is:
\begin{equation}
  \bfA_S[i,j] = \bfA[i,j] \cdot \sigmoid\!\bigl(\bfM_A[i,j]\bigr),
  \quad (i,j) \in \calE_c^v.
  \label{eq:edge_mask}
\end{equation}
This soft-masked adjacency replaces the binary one in all GNN message-passing
computations, maintaining differentiability. In practice, the objective in
Eq.~\eqref{eq:class_specific} is evaluated by a single forward pass through
the soft-weighted graph. This is a mean-field approximation, discussed in
Section~\ref{sec:limitations}.

\paragraph{Feature mask.}
\citet{ying2019gnnexplainer} intend the feature mask $F \in \{0,1\}^d$ as a
\emph{binary} selector, $\bfx_j^F = \bfx_j \odot F$. They warn against
naively relaxing $F$ to a continuous $\sigmoid(\bfm_F) \in (0,1)^d$
multiplied directly into $\bfx_j$, as done for the edge mask
(Eq.~\ref{eq:edge_mask}): an important feature whose value is already close
to zero becomes indistinguishable from an unimportant one, since scaling it
near zero either way gives almost no gradient signal about its true
importance. Their proposed fix, a marginal-distribution reparametrization
trick, is summarized in Appendix~\ref{app:featmask}.

The PyTorch Geometric implementation used in our experiments
(\texttt{node\_mask\_type='attributes'}) skips this reparametrization entirely.
It learns an unconstrained mask $\bfM_F \in \reals^{n \times d}$ (one entry per
node \emph{and} feature, not shared across nodes as originally proposed) and
applies it exactly like the edge mask,
\begin{equation}
  \bfX_S^F = \bfX_c^v \odot \sigmoid(\bfM_F),
  \label{eq:feat_mask}
\end{equation}
i.e., the same near-zero-value blind spot is present in the library. A
practitioner explaining informative small-valued features should be aware
\GNNEx may under-report their importance; every node feature in our
experiments is a $0/1$ indicator, however, so this case does not arise and
Eq.~\eqref{eq:feat_mask} is adequate for our purposes.

\subsection{From a Discrete Search to a Single Forward Pass: A Flawed Step}
\label{subsec:jensen}

Eqs.~\eqref{eq:cond_ent}--\eqref{eq:class_specific} still require searching over
discrete subgraphs $\calG_S$. \citet{ying2019gnnexplainer} make this tractable
by treating $\calG_S \sim \mathcal{Q}$ as a random variable drawn from a
distribution $\mathcal{Q}$ over subgraphs, and by minimizing
$\mathbb{E}_{\calG_S \sim \mathcal{Q}}\bigl[H(Y \mid \calG_S, \bfX_S^F)\bigr]$.
They then invoke \emph{Jensen's inequality} which
states that for a convex\footnote{A function $\varphi$ is \emph{convex} when
$\varphi(\lambda a + (1{-}\lambda)b) \leq \lambda\varphi(a) + (1{-}\lambda)\varphi(b)$
for all $a,b$ in its domain and $\lambda\in[0,1]$.} $\varphi$ and a random variable $Z$ taking values in
its domain, $\varphi(\mathbb{E}[Z]) \leq \mathbb{E}[\varphi(Z)]$. Here
$\varphi = H(Y\mid\cdot,\bfX_S^F)$ is regarded as a function of the fractional
adjacency matrix $\bfA_S$ of Eq.~\eqref{eq:edge_mask}. Ying et al.\ use this
to replace the expectation above with
$H\bigl(Y \mid \mathbb{E}_{\mathcal{Q}}[\calG_S], \bfX_S^F\bigr)$, licensing
the single soft-masked forward pass of Eq.~\eqref{eq:edge_mask} in place of
an intractable average; they state that, ``with convexity assumption,''
this substitution is an \emph{upper} bound on the original expectation.

This is backwards. Under the very convexity assumption invoked (that
$H(Y \mid \calG_S, \cdot)$ is convex in $\calG_S$), Jensen's inequality actually
gives $H(Y \mid \mathbb{E}_{\mathcal{Q}}[\calG_S], \bfX_S^F) \leq
\mathbb{E}_{\mathcal{Q}}\bigl[H(Y \mid \calG_S, \bfX_S^F)\bigr]$: a \emph{lower}
bound, not an upper one, and minimizing a lower bound gives no guarantee about
the quantity of actual interest. Worse, the authors immediately note that
``due to the complexity of neural networks, the convexity assumption does not
hold'' for the models \GNNEx targets, so even the mislabeled bound lacks a
valid premise. This step has no rigorous justification; we adopt it only
because it is what every practical implementation optimizes, and because
its value rests on strong empirical results (Section~\ref{sec:results}),
not on the derivation.

\subsection{Optimization Objective}

Substituting the mask-parameterized adjacency (Eq.~\ref{eq:edge_mask}) and
features (Eq.~\ref{eq:feat_mask}) into the class-specific formulation
(Eq.~\ref{eq:class_specific}), the full objective adds sparsity and
discreteness regularization:
\begin{align}
  \min_{\bfM_A,\, \bfM_F}\;
  -\log P_\Phi(\hat{y}_v \mid \bfA_S, \bfX_S^F)
  &+ \alpha_1 \!\sum_{(i,j)}\! \sigmoid(\bfM_A[i,j])
   + \alpha_2 \!\sum_{(i,j)}\! H\!\bigl(\sigmoid(\bfM_A[i,j])\bigr)
  \nonumber\\
  &+ \alpha_3 \sum_{v,k} \sigmoid(\bfM_F[v,k])
   + \alpha_4 \sum_{v,k} H\!\bigl(\sigmoid(\bfM_F[v,k])\bigr),
  \label{eq:full_obj}
\end{align}
where $H(p) = -p\log p - (1-p)\log(1-p)$. The $\ell_1$ terms ($\alpha_1,
\alpha_3$) enforce sparsity; the binary-entropy terms ($\alpha_2, \alpha_4$)
push mask values toward 0 or 1. The final subgraph is obtained by thresholding
the edge mask at 0.5 or retaining the top-$K$ edges.

\subsection{Local Optima in the Mask Optimization}
\label{subsec:localopt}

Because \GNNEx solves Eq.~\eqref{eq:full_obj} by gradient descent, explanation
quality depends on which local optimum the optimizer reaches.
\citet{ying2019gnnexplainer} already anticipate this: since their objective
is non-convex (Section~\ref{subsec:jensen}), gradient descent finds only
\emph{some} local minimum, which they report empirically tends to be a
high-quality explanation. This is therefore not a new discovery; our
contribution is a concrete investigation of when it instead settles on a
poor one, particularly when every strict sub-structure of the correct
explanation already maintains reasonable \emph{prediction confidence}, the
probability assigned to the predicted class.

Consider a node $v$ in a $k$-clique in a sparse graph, where the GNN assigns
high confidence to class~1 whenever the local neighborhood is dense. A
triangle (3-clique) embedded within the $k$-clique is already denser than
the Erd\H{o}s-R\'{e}nyi \citep{erdos1959random} background, so the triangle
alone may already preserve confident class-1 prediction. The gradient of
Eq.~\eqref{eq:full_obj} is then small: removing a triangle edge would
decrease confidence, and the entropy regularization discourages adding
more. The optimizer stops at this \emph{sub-clique local minimum} even
though the true explanation is the full $k$-clique, and as $k$ grows, the
number of competing sub-clique local optima grows exponentially. A weak
classifier compounds this: if the GNN does not reliably predict clique
nodes, the gradient signal is small throughout, and convergence to any
meaningful explanation becomes unlikely.

\subsection{Anchor Features and Unfaithful Explanations in \GNNEx}
\label{subsec:anchor}

The unfaithfulness result reviewed in Section~\ref{sec:related} concerns
SE-GNNs, where $e$ and $g$ are optimized jointly. \GNNEx has no such
freedom, so we state the mechanism precisely for our setting, to fix what
\emph{would} constitute the same failure here and motivate a testable,
weaker analogue.

\citet{azzolin2026gnnexplanations} study graph classification, where anchor
\emph{nodes} recur across every graph. A set of anchor nodes need not lie
inside a given node's computation graph $\calG_c^v$ (Section~\ref{sec:background}),
so it cannot serve as a node-level anchor in general, unlike a node
\emph{feature}, which is present at every node. This is why we use color,
not designated nodes, below.

Augment the multi-size clique dataset with two constant \emph{color}
features per node, $\text{color}_1 \equiv 1$ and $\text{color}_2 \equiv 2$,
taking the exact same value on \emph{every} node. Both are anchor features
(Section~\ref{sec:related}): $\text{color}_1 = 1$ and $\text{color}_2 = 2$
hold for every node $v \in \calV$ regardless of its true label
$y_v \in \{0,1\}$, so neither carries any information about clique
membership. A hypothetical SE-GNN pair $(e,g)$ could still reach perfect
accuracy while reporting only one of these two, always-true facts as its
explanation. Let $e(v) := \{v,\text{color}_1\}$ whenever $v$ truly belongs
to a clique, determined by $e$ examining the \emph{entire} graph internally,
and $e(v) := \{v,\text{color}_2\}$ otherwise. With
$g(\text{color}_1)=\text{clique}$ and $g(\text{color}_2)=\text{non-clique}$,
$f=g\circ e$ gives $f(G) = g(e(G)) = g(\text{color}_1) = \text{clique}$ for
every clique node. This is an instance of Theorem~1 of
\citet{azzolin2026gnnexplanations}: $f$ is perfectly accurate, yet its
reported explanation, identically true for every node in the graph, is
entirely unrelated to the true decision-relevant structure, visible only
inside $e$'s hidden computation and never revealed.

\GNNEx cannot build this encoding, since $\Phi$ is fixed and there is no $e$
to optimize: this is exactly what separates a post-hoc explainer from an
SE-GNN. In Theorem~1, the failure lives in the \emph{joint training} of
$(e,g)$, so it occurs regardless of how $\Phi$ treats the anchor set; a
post-hoc explainer can only report what an already-trained $\Phi$ does. Any
anchor-style failure here must therefore already exist inside $\Phi$
itself, whether the anchor is a set of graph-level nodes
\citep{azzolin2026gnnexplanations} or our node-level features: an
ideally-trained $\Phi$ would ignore them, and \GNNEx would then ignore them
too. If $\Phi$'s training instead assigns any non-zero, spuriously small
weight to $\text{color}_1$ or $\text{color}_2$, \GNNEx's feature mask has a
real gradient to exploit, and the mean-field optimizer, already prone to
local optima (Section~\ref{subsec:localopt}), may inflate that weight into
an unfaithful explanation. We test this next.


\section{Experiments}
\label{sec:experiments}

\subsection{Synthetic Dataset}

\paragraph{Multi-Size Clique Detection with a Color Anchor.}
To study local optima, we embed cliques of sizes 3, 4, and 5 in a single sparse
Erd\H{o}s-R\'{e}nyi graph \citep{erdos1959random} with $n=200$ nodes, each
pair independently connected with probability $p=0.02$: 8 triangles,
6 four-cliques, and 5 five-cliques, covering 73 of the 200 nodes. Labels are
binary (0 = non-clique, 1 = clique node). Node features are two constant
\emph{color} features, $\text{color}_1 \equiv 1$ and $\text{color}_2 \equiv
2$, with the same value on every node. These are the anchor features of
Section~\ref{subsec:anchor}, uninformative about the label by construction.
The ground-truth explanation for a node in a $k$-clique consists of the
$\binom{k}{2}$ undirected edges of that clique.

This setting exposes the local-optima behavior of Section~\ref{subsec:localopt}:
a triangle embedded in a 5-clique is already a locally consistent
explanation, yet represents only 3 of the 10 true explanatory edges.

\subsection{GNN Architecture and Training}

We train a 3-layer GCN \citep{kipf2017semi} with 64 hidden units, ReLU
activations, and dropout ($p=0.5$), for 300 epochs using Adam
\citep{kingma2015adam} with learning rate $10^{-2}$ and weight decay
$5\times10^{-4}$. We apply class-weighted cross-entropy
($w_1 = n_{\mathrm{base}}/n_{\mathrm{clique}}$) to compensate for the 127:73
class imbalance.

\subsection{Explainer Configuration and Evaluation Metrics}

We use the \GNNEx implementation in PyTorch Geometric \citep{fey2019fast},
run for 150 gradient-descent steps per node at learning rate $10^{-2}$, learning
both edge and node-feature attribute masks.

\paragraph{Classification accuracy.}
We first verify the GNN itself by reporting \emph{accuracy}: the fraction of
correctly classified nodes, $\mathrm{Acc} = |\{v : \hat{y}_v = y_v\}| / |\calV|$.

\paragraph{Explanation quality.}
Edges with mask weight $\sigmoid(\bfM_A[i,j]) > 0.5$ are \emph{selected}; the
rest are discarded. Let $S$ be the selected edge set and $T$ the ground-truth
edge set for a given node. We report precision, recall, and their harmonic
mean F1:
\begin{equation}
  P = \frac{|S \cap T|}{|S|}, \qquad
  R = \frac{|S \cap T|}{|T|}, \qquad
  F_1 = \frac{2PR}{P + R}.
  \label{eq:f1}
\end{equation}
\citet{azzolin2026gnnexplanations} review alternative faithfulness metrics
that probe an explanation by perturbing the graph rather than comparing to
ground truth (e.g., Fidelity, or their own EST metric), and
\citet{faber2021when} and \citet{agarwal2022probing} argue such comparison
can itself be misleading (Section~\ref{sec:limitations}). We use F1 against
the known planted ground truth anyway, as it remains the most widely
reported metric on synthetic benchmarks with known explanations.

We additionally report the \emph{explanation clique size}: the largest clique
found among the true clique nodes in the thresholded explanation subgraph
(via \texttt{networkx.find\_cliques}).

\paragraph{Anchor-feature (color) weight.}
\GNNEx optimizes a separate feature mask for each explained node; write
$\bfM_F^{(v)}$ for the mask learned while explaining $v$
(Eq.~\ref{eq:feat_mask}), and
$\bfw_v := \sigmoid\bigl(\bfM_F^{(v)}[v,:]\bigr) \in (0,1)^2$ for $v$'s own
row of it, the weight \GNNEx assigns to $v$'s own $\text{color}_1$ and
$\text{color}_2$ entries. A faithful explainer should assign $\bfw_v$ a low
value that does not depend on $v$'s class. A class-dependent asymmetry,
e.g., a high $\text{color}_1$ weight on clique nodes paired with a high
$\text{color}_2$ weight on non-clique nodes, would be a soft version of the
deterministic encoding constructed in Section~\ref{subsec:anchor}, though
Section~\ref{sec:limitations} argues this would indict $\Phi$'s training,
not \GNNEx.


\section{Results and Discussion}
\label{sec:results}

\subsection{Multi-Size Clique Detection}

The GCN reaches 71.5\% overall accuracy: 90.4\% on clique nodes but only
60.6\% on base nodes, so the class-weighted loss favors confidently
detecting cliques at the cost of base-node accuracy, which directly affects
the explainer.

Table~\ref{tab:results} shows a sharp degradation with clique size: F1
falls from 0.124 ($k=3$) to 0.000 for $k=4$ and $k=5$, where the optimizer
recovers no edge within the true clique. This confirms the local-optima
mechanism of Section~\ref{subsec:localopt}: for $k=3$ the target has only 3
edges, so the optimizer sometimes reaches it before converging, while for
larger cliques any single clique edge is already a locally sufficient
explanation.

\begin{table}[t]
  \centering
  \caption{Explanation accuracy on the clique benchmark.
    Edges are thresholded at mask weight $>0.5$.
    F1 values are mean $\pm$ std over all explained nodes.
    ``Expl.\ clique'' reports the average size of the largest clique
    recovered among the true clique nodes in the explanation subgraph.}
  \label{tab:results}
  \smallskip
  \begin{tabular}{lcccc}
    \toprule
    Clique size & Precision & Recall & F1 & Expl.\ clique \\
    \midrule
    $k=3$ (triangles)  & 0.174 & 0.097 & $0.124 \pm 0.204$ & 1.29 / 3 \\
    $k=4$              & 0.000 & 0.000 & $0.000 \pm 0.000$ & 1.00 / 4 \\
    $k=5$              & 0.000 & 0.000 & $0.000 \pm 0.000$ & 1.00 / 5 \\
    All clique nodes   & 0.073 & 0.041 & $0.052 \pm 0.146$ & -- \\
    \bottomrule
  \end{tabular}
\end{table}

\subsection{Anchor-Feature (Color) Results}
\label{subsec:anchor_results}

We measure the mean color-feature weight $\bfw_v$ (Section~\ref{sec:experiments})
over 57 explained clique nodes and 50 explained non-clique nodes. A faithful
explainer would assign both $\text{color}_1$ and $\text{color}_2$ a weight
close to 0, equally for both classes, since both are constant and carry no
information about $y_v$. Neither is suppressed to near zero: clique nodes
receive $\text{color}_1=0.186$ and $\text{color}_2=0.266$, and non-clique
nodes receive $\text{color}_1=0.188$ and $\text{color}_2=0.235$, so the
default regularization does not drive a genuinely irrelevant feature to 0
within 150 steps.

The class-dependent gap, however, is small: $\text{color}_1$ shows
essentially none (0.186 vs.\ 0.188), and $\text{color}_2$ shows only a
modest one (0.266 vs.\ 0.235) that does not track classifier confidence as a
compensation account would predict, since clique nodes are now the
\emph{easier} class (90.4\% vs.\ 60.6\% accuracy) yet receive the higher
weight. \GNNEx is thus only weakly susceptible to the
Section~\ref{subsec:anchor} failure once the anchor feature is made truly
constant: it does not fabricate a clean, deterministic encoding
($g(\text{color}_1){=}\text{clique}$), but the small, unexplained gap on
$\text{color}_2$ shows it is not perfectly faithful either, a gap
Section~\ref{sec:limitations} attributes to $\Phi$'s training rather than
\GNNEx.


\section{Limitations}
\label{sec:limitations}

Three further limitations qualify our results. First, the mean-field
relaxation (Section~\ref{subsec:jensen}) evaluates a soft-masked graph
outside the binary-adjacency distribution $\Phi$ was trained on: $\Phi$
never saw an edge held at half-strength, so its behavior there is
unconstrained and may not match how it would behave if that edge were
genuinely removed. The explanation may therefore reflect an effective model
that differs from $\Phi$ \citep{yuan2022explainability};
\citet{agarwal2022probing} show this introduces a bounded but non-zero
faithfulness gap. Second, each \GNNEx call is an independent optimization
with no principled way to set $\alpha_1,\ldots,\alpha_4$
(Eq.~\ref{eq:full_obj}), making it less stable across structurally
identical nodes than amortized methods like PGExplainer
\citep{luo2020parameterized}.

Third, \citet{faber2021when} and \citet{agarwal2022probing} argue that
comparing an explanation to ground truth can itself be misleading: a
badly-trained GNN may not use the planted motif at all, so a low score
reflects the classifier, not the explainer. Our clique GCN reaches 90.4\%
accuracy on clique nodes but only 60.6\% on base nodes, so some of the F1
gap may not be \GNNEx's fault. The same applies even more sharply to our
anchor-feature check (Section~\ref{subsec:anchor_results}): since \GNNEx
only reports what $\Phi$ already relies on (Section~\ref{subsec:anchor}),
any non-zero color weight we measure is a statement about $\Phi$'s
training, not \GNNEx. We report it anyway, for simplicity and because
evaluating faithfulness metrics is not this report's focus.


\section{Conclusion and Future Work}
\label{sec:conclusion}

We have presented a detailed study of \GNNEx. Along the way we identified a
genuine error in the original derivation (Section~\ref{subsec:jensen}): the
step that replaces an intractable expectation over subgraphs with a single
soft-masked forward pass invokes Jensen's inequality in the wrong
direction, so the continuous relaxation underlying \GNNEx is best
understood as a heuristic, not a provable bound. Our main structural
contribution concretely investigates the local-optima behavior the
original authors already anticipate but do not characterize, confirmed on
the clique benchmark: F1 drops from 0.12 for triangles to 0.00 for
5-cliques. Our main empirical contribution adapts a 2026 unfaithfulness
result for self-explainable GNNs into a testable question for \GNNEx: once
the anchor features are made truly constant, its feature mask cannot
reproduce the clean, deterministic failure of a jointly-trained extractor,
and only a weak, partly unexplained trace survives.

Future work could address local optima with a multi-restart strategy or a
coarse combinatorial initialization, and revisit the mean-field relaxation
with a correctly-directed bound. The color-anchor construction offers a
controllable way to probe post-hoc explainers for a failure mode so far
studied mainly for self-explainable models.

\bibliographystyle{plainnat}
\bibliography{references}

\appendix

\section{Implementation Details}
\label{app:impl}

All experiments use Python~3.13, PyTorch~2.12, and PyTorch Geometric~2.8
\citep{fey2019fast}. The full source code is provided as
\texttt{gnnexplainer\_synthetic.py}. Random seeds are fixed throughout for
reproducibility. The \GNNEx implementation in PyTorch Geometric requires
specifying the explanation type (\texttt{`model'}), mask types
(\texttt{edge\_mask\_type='object'} and \texttt{node\_mask\_type='attributes'}),
and model configuration. Using \texttt{mode='multiclass\_classification'} with
two output classes handles binary classification without requiring changes to the
GCN output layer. As verified against the library source
(\texttt{torch\_geometric/explain/algorithm/gnn\_explainer.py}), the
\texttt{'attributes'} mask type creates one learnable weight per node \emph{and}
feature (Eq.~\ref{eq:feat_mask}), which is what we read out per node to obtain
the color-anchor weights of Section~\ref{subsec:anchor_results}.

\section{Feature-Mask Reparametrization: The Original Proposal}
\label{app:featmask}

\citet{ying2019gnnexplainer} address the near-zero-value blind spot of
Section~\ref{sec:method} by not multiplying discarded feature entries by zero
(or a value near it) at all. Instead, they treat the masked feature as a random
variable and replace discarded entries with samples from that feature's own
empirical marginal distribution, so a discarded feature is swapped for a
\emph{typical} value rather than suppressed toward zero. This avoids the
confound where "small mask weight" and "small feature value" look identical to
the optimizer. Concretely, using a reparametrization trick to keep the sampling
step differentiable,
\begin{equation}
  \bfx_j^F = \bfz + (\bfx_j - \bfz) \odot F,
  \qquad \textstyle\sum_k F_k \leq K_F,
  \label{eq:feat_reparam}
\end{equation}
where $\bfz$ is sampled from the empirical marginal distribution of the node
features and $K_F$ caps the number of features kept in the explanation; the
expectation over $\bfz$ is estimated by Monte Carlo sampling during training.
Eq.~\eqref{eq:feat_reparam} is the mechanism referenced in
Section~\ref{sec:method}: PyTorch Geometric's \GNNEx implementation does not
carry it out, using the plain sigmoid mask of Eq.~\eqref{eq:feat_mask} instead.

\section{The RBGV Example}
\label{app:rbgv}

\citet{azzolin2026gnnexplanations} illustrate the anchor-set construction of
Section~\ref{sec:related} with a synthetic graph-classification dataset, RBGV:
every graph contains red and blue nodes, randomly connected, whose relative
counts ($\#\text{blue}$ vs.\ $\#\text{red}$) determine a binary label, plus one
isolated green node and one isolated violet node present in \emph{every}
instance, forming an anchor set. An SE-GNN can reach perfect task accuracy
while reporting only the anchor nodes as its explanation: let the extractor
output the green node when reds dominate and the violet node otherwise,
\begin{equation}
  e(G) = \begin{cases} u_{\text{green}} & \text{if } \#\text{red} \geq \#\text{blue} \\
                        u_{\text{violet}} & \text{if } \#\text{blue} > \#\text{red}, \end{cases}
  \qquad
  g(R) = \begin{cases} 0 & \text{if } R = u_{\text{green}} \\
                        1 & \text{if } R = u_{\text{violet}}. \end{cases}
\end{equation}
The composition $f = g \circ e$ is perfectly accurate, yet its reported
explanation is color-blind to red and blue, the only nodes the label
actually depends on. This is because $e$ must internally count red and blue
nodes to decide which anchor to report, but never reveals that computation.
This is the template Section~\ref{subsec:anchor} adapts to our clique
dataset.

\end{document}
